% Template de TCC — Curso de Graduação em Geologia — UFSC
% Compile com XeLaTeX (Overleaf) ou Tectonic.
\documentclass[
  12pt,
  oneside,
  a4paper,
  chapter=TITLE,
  section=TITLE,
  english,
  brazil
]{abntex2}

% -----------------------------------------------------------------------------
% DADOS DO TRABALHO — edite somente este arquivo para preencher os metadados.
% -----------------------------------------------------------------------------

\newcommand{\AutorTCC}{Maria Eulália Alberton}
\newcommand{\TituloTCC}{Caracterização de depósitos de movimento de massa na porção norte da Bacia de Santos}
\newcommand{\SubtituloTCC}{} % deixe vazio se não houver subtítulo

\newcommand{\OrientadorTCC}{Prof. Nome do orientador, Dr.}
\newcommand{\CoorientadorTCC}{} % deixe vazio se não houver coorientação

\newcommand{\MembroUmTCC}{Prof. Nome do primeiro membro, Dr.}
\newcommand{\InstituicaoMembroUm}{Universidade Federal de Santa Catarina}
\newcommand{\MembroDoisTCC}{Prof. Nome do segundo membro, Dr.}
\newcommand{\InstituicaoMembroDois}{Universidade Federal de Santa Catarina}

\newcommand{\CoordenadorTCC}{Prof. Nome da coordenação, Dr.}
\newcommand{\DataDefesaTCC}{dia de mês de ano}
\newcommand{\CidadeTCC}{Florianópolis}
\newcommand{\AnoTCC}{2026}

\newcommand{\InstituicaoTCC}{Universidade Federal de Santa Catarina}
\newcommand{\CentroTCC}{Centro de Filosofia e Ciências Humanas}
\newcommand{\DepartamentoTCC}{Departamento de Geologia}
\newcommand{\CursoTCC}{Curso de Graduação em Geologia}
\newcommand{\GrauTCC}{Bacharel em Geologia}

\newcommand{\PreambuloTCC}{Trabalho de Conclusão de Curso submetido ao Curso de Graduação em Geologia do Centro de Filosofia e Ciências Humanas da Universidade Federal de Santa Catarina como requisito parcial para a obtenção do título de Bacharel em Geologia.}

% Dados usados nos metadados do PDF e pelos mecanismos do abnTeX2.
\autor{\AutorTCC}
\titulo{\TituloTCC}
\local{\CidadeTCC}
\data{\AnoTCC}
\orientador{\OrientadorTCC}
\preambulo{\PreambuloTCC}

\usepackage{estilo/tcc-geologia-ufsc}

\begin{document}
\selectlanguage{brazil}
\frenchspacing
\OnehalfSpacing

% -----------------------------------------------------------------------------
% ELEMENTOS PRÉ-TEXTUAIS
% -----------------------------------------------------------------------------
\pretextual
\imprimircapaGeologia
\imprimirfolhaderostoGeologia
\imprimirfichacatalograficaModelo
\imprimirfolhadeaprovacaoGeologia
\input{pre-textuais/dedicatoria}
\input{pre-textuais/agradecimentos}
\input{pre-textuais/epigrafe}
\input{pre-textuais/resumo}
\input{pre-textuais/abstract}

\pdfbookmark[0]{\listfigurename}{lof}
\listoffigures*
\cleardoublepage

\pdfbookmark[0]{\listtablename}{lot}
\listoftables*
\cleardoublepage

\pdfbookmark[0]{\listofquadrosname}{loq}
\listofquadros
\cleardoublepage

\input{pre-textuais/siglas}
\input{pre-textuais/simbolos}

\pdfbookmark[0]{\contentsname}{toc}
\tableofcontents*
\cleardoublepage

% -----------------------------------------------------------------------------
% ELEMENTOS TEXTUAIS
% Em abnTeX2, \chapter representa uma seção primária na apresentação ABNT.
% -----------------------------------------------------------------------------
\textual
\pagestyle{tccgeologia}
\chapter{Introdução}\label{cap:introducao}

\TextoModelo{Apresente o tema, delimite o problema geológico, sintetize o estado do conhecimento, justifique o estudo e antecipe a abordagem utilizada. A introdução deve conduzir naturalmente aos objetivos.}

Como exemplo de citação no sistema autor-data, uma referência pode ser integrada ao texto por meio de \citeonline{hasui2012} ou apresentada entre parênteses \cite{press2006}. Substitua estas referências pelas fontes efetivamente utilizadas.

\section{Objetivos}\label{sec:objetivos}

\subsection{Objetivo geral}

\TextoModelo{Declare o resultado científico principal que o trabalho pretende alcançar.}

\subsection{Objetivos específicos}

\begin{itemize}
  \item \TextoModelo{caracterizar as unidades geológicas da área de estudo;}
  \item \TextoModelo{descrever os procedimentos de campo, laboratório e gabinete;}
  \item \TextoModelo{integrar e interpretar os dados obtidos;}
  \item \TextoModelo{discutir as implicações geológicas dos resultados.}
\end{itemize}

\section{Justificativa}

Apesar de serem menos estudados quando comparado aos deslizamentos terrestres, os movimentos de massa em ambiente marinho também são importantes em termos sociais, econômicos e ecológicos (Clare et al., 2018), uma vez que são fontes potenciais a gerar danos (Vanneste et al., 2014). Frequentemente, esses deslizamentos submarinos geram fluxos de dezenas a centenas de quilômetros, capazes de danificar a infraestrutura estratégica no fundo do mar, como cabos de telecomunicações, plataformas de produção e dutos de hidrocarbonetos (Pope et al., 2016), os quais estão presentes na região norte da Bacia de Santos.

Sob soterramento, os depósitos de deslizamentos subaquáticos são elementos importantes nos sistemas petrolíferos, pois influenciam a distribuição dos reservatórios (Kneller et al., 2016), podem atuar como selos (Cardona et al., 2016) e também como potenciais reservatórios (Meckel, 2011).

Um dos maiores desafios de estudos sobre transportes de massa submarinos é a sua natureza, em mar aberto e longe da costa, o que dificulta a aquisição dos dados bem como a falta de informação consistente, por vezes, de baixa resolução (Clare et al., 2018). Assim, convergir conhecimentos da atividade sismoestratigráfica e sismológica, visando a identificação e caracterização de movimentos de massa, dará suporte para avançar no entendimento de processos atuantes na margem continental sudeste e, consequentemente, no planejamento de ações da indústria do petróleo na implantação e monitoramento de infraestrutura de exploração e produção de hidrocarbonetos.

\chapter{Fundamentação teórica}\label{cap:fundamentacao}

\section{Sísmica de reflexão}

A sísmica de reflexão é um dos métodos geofísicos mais efetivos para mapear a geologia de uma determinada área, principalmente em ambientes submersos, utilizando fontes artificiais. Esse método baseia-se no princípio de que, quando um pulso sísmico é provocado próximo à superfície da coluna d'água por uma fonte artificial, ondas mecânicas se propagam em direção ao interior da Terra até encontrarem interfaces que delimitam camadas litológicas com diferentes impedâncias acústicas, nas quais parte de sua energia é refletida (Tóth, 2011).

A impedância acústica ($Z$) é uma propriedade que pode ser determinada pelo produto entre a densidade ($\rho$) do meio e a velocidade de propagação ($c$) da onda, considerando-se uma onda plana com incidência normal à interface. Assim, a impedância acústica de um meio pode ser expressa pela relação $Z = \rho c$ (Yilmaz, 2001). O fenômeno físico da reflexão das ondas acústicas ocorre quando existem contrastes representativos de impedância entre os meios atravessados pelo sinal emitido (Tóth, 2011). Nessas interfaces, parte da energia da onda retorna em direção à superfície, enquanto outra parte é transmitida e continua sua propagação até encontrar uma nova camada com impedância acústica distinta, repetindo-se sucessivamente os processos de reflexão e transmissão.

A técnica baseia-se na análise dos intervalos de tempo decorridos entre a emissão dos sinais pela fonte artificial e a chegada das ondas refletidas aos receptores, bem como de suas amplitudes. Esse intervalo é denominado tempo duplo de percurso, ou \textit{two-way travel time} (TWT), pois corresponde ao tempo necessário para que a onda se desloque da fonte até uma determinada interface e retorne à superfície. Dessa forma, os refletores presentes nas seções sísmicas são inicialmente posicionados em função do tempo, geralmente expresso em segundos ou milissegundos, e não diretamente em profundidade. A conversão do tempo em profundidade depende do conhecimento ou da estimativa das velocidades de propagação das ondas nos diferentes materiais atravessados (Yilmaz, 2001).

Cada receptor registra, ao longo do tempo, as variações de amplitude das ondas que retornam à superfície. Esse registro individual é denominado traço sísmico e representa a resposta do meio geológico ao pulso emitido em uma determinada posição. A organização de vários traços sísmicos lado a lado forma uma seção sísmica, na qual a continuidade, a geometria e a amplitude das reflexões permitem identificar e interpretar interfaces e estruturas geológicas em subsuperfície (Kearey, 2009; Yilmaz, 2001). Assim, o método de sísmica de reflexão fundamenta-se no processamento e na posterior interpretação das reflexões registradas na superfície (Kearey, 2009).

\subsection{Cosseno de fase instantânea}

O atributo sísmico cosseno da fase instantânea é obtido a partir do cálculo do cosseno da fase instantânea do sinal sísmico, conforme a Equação~\ref{eq:cosseno-fase}.

\begin{equation}
  C(t) = \cos\bigl(\phi(t)\bigr)
  \label{eq:cosseno-fase}
\end{equation}

em que $C(t)$ corresponde ao cosseno da fase instantânea e $\phi(t)$ representa a fase instantânea do sinal sísmico no tempo $t$. A fase instantânea é determinada a partir do sinal analítico, cuja parte imaginária é obtida por meio da transformada de Hilbert. Esse procedimento permite separar as informações de fase e amplitude do traço sísmico (Chopra \& Marfurt, 2007). Consequentemente, o cálculo do cosseno da fase torna-se independente dos valores de amplitude, promovendo uma normalização do traço sísmico e removendo as variações relacionadas à intensidade do sinal (Meneses, 2010).

Ao reduzir a influência dos contrastes de amplitude, esse atributo evidencia a geometria e a configuração dos refletores, permitindo que feições anteriormente pouco perceptíveis sejam visualizadas com maior clareza (Chopra \& Marfurt, 2007). Dessa forma, o cosseno da fase instantânea contribui para realçar a continuidade lateral dos refletores, as mudanças no padrão de acamamento e os limites entre sequências sísmicas, auxiliando a interpretação estratigráfica e estrutural dos dados (Alvarenga, 2016).

\section{Contexto geológico regional}

\TextoModelo{Apresente as unidades litoestratigráficas, o contexto tectônico, a evolução geológica e os trabalhos anteriores necessários à compreensão da área. Organize a revisão do regional para o local.}

\section{Geologia local}

\TextoModelo{Descreva as unidades mapeadas, suas relações de contato, estruturas, metamorfismo, alteração, geomorfologia e demais aspectos pertinentes ao tema.}

\chapter{Área de estudo}\label{cap:area-estudo}

% Conteúdo a ser desenvolvido no Google Docs e transferido para esta seção.

\chapter{Materiais e métodos}\label{cap:metodos}

\section{Pesquisa bibliográfica e bases de dados}

\TextoModelo{Indique bases consultadas, critérios de seleção, mapas, imagens, modelos digitais de elevação, dados geoquímicos e demais produtos utilizados.}

\section{Etapa de campo}

\TextoModelo{Informe período, escala de trabalho, equipamentos, critérios de descrição de afloramentos, amostragem, controle de qualidade e registro de coordenadas.}

\begin{quadro}
\caption{Exemplo de organização do levantamento de campo}\label{qua:campo}
\centering
\begin{tabularx}{\textwidth}{>{\bfseries}l X l}
\toprule
Atividade & Registro mínimo & Produto \\
\midrule
Descrição de afloramento & Litologia, textura, estruturas, alteração e fotografias & Ficha de ponto \\
Medida estrutural & Tipo de estrutura, atitude, método e qualidade & Banco estrutural \\
Amostragem & Código, material, finalidade e posição & Cadeia de custódia \\
\bottomrule
\end{tabularx}
\FonteTCC{Elaboração própria}
\end{quadro}

\section{Etapa laboratorial}

\TextoModelo{Descreva preparação das amostras, métodos analíticos, equipamentos, padrões, limites de detecção, brancos, duplicatas e tratamento de incertezas.}

\section{Tratamento e integração dos dados}

\TextoModelo{Explique os procedimentos estatísticos, petrográficos, geoquímicos, geocronológicos, geofísicos, geotécnicos ou de geoprocessamento adotados. Declare softwares e versões quando forem relevantes à reprodutibilidade.}

\chapter{Resultados}\label{cap:resultados}

\TextoModelo{Apresente as observações e medições de forma organizada, evitando antecipar toda a interpretação. Use a mesma ordem lógica dos métodos ou dos objetivos específicos.}

\section{Unidades geológicas}

\TextoModelo{Descreva os resultados de mapeamento, petrografia, estratigrafia ou caracterização de materiais. Diferencie observação de interpretação.}

\begin{table}[H]
\centering
\caption{Exemplo de síntese das unidades reconhecidas}\label{tab:unidades}
\begin{tabularx}{\textwidth}{l X X}
\toprule
Unidade & Características de campo & Interpretação preliminar \\
\midrule
U1 & \TextoModelo{litologia, textura, estruturas e contatos} & \TextoModelo{ambiente ou processo} \\
U2 & \TextoModelo{litologia, textura, estruturas e contatos} & \TextoModelo{ambiente ou processo} \\
\bottomrule
\end{tabularx}
\FonteTCC{Elaboração própria}
\end{table}

\section{Dados estruturais}

\TextoModelo{Apresente famílias de estruturas, estatísticas descritivas, estereogramas, relações de corte e indicadores cinemáticos.}

\section{Resultados analíticos}

\TextoModelo{Apresente os resultados laboratoriais com unidades, precisão, limites analíticos e identificação inequívoca das amostras.}

\figuraplaceholder{Perfil geológico esquemático da área de estudo}{Elaboração própria}{fig:perfil-geologico}

\chapter{Discussão}\label{cap:discussao}

\TextoModelo{Interprete os resultados em relação aos objetivos, confronte-os com a literatura, avalie hipóteses alternativas e explicite limitações. Evite apenas repetir a seção anterior.}

\section{Integração dos resultados}

\TextoModelo{Relacione dados de campo, laboratório e gabinete. Mostre quais evidências sustentam cada interpretação.}

\section{Modelo evolutivo ou conceitual}

\TextoModelo{Quando pertinente, proponha uma sequência de eventos, um modelo deposicional, estrutural, petrológico, hidrogeológico, geotécnico ou ambiental.}

\section{Limitações e incertezas}

\TextoModelo{Discuta cobertura espacial, representatividade das amostras, resolução, erros analíticos, incertezas cartográficas e hipóteses que condicionam as conclusões.}

\chapter{Conclusões}\label{cap:conclusoes}

\TextoModelo{Responda diretamente aos objetivos. Sintetize as conclusões sustentadas pelos resultados, sem introduzir dados novos. Indique recomendações e trabalhos futuros apenas quando decorrerem do estudo.}

\begin{itemize}
  \item \TextoModelo{conclusão principal associada ao objetivo geral;}
  \item \TextoModelo{conclusão associada ao primeiro objetivo específico;}
  \item \TextoModelo{conclusão associada ao segundo objetivo específico;}
  \item \TextoModelo{recomendação ou questão em aberto.}
\end{itemize}


% -----------------------------------------------------------------------------
% ELEMENTOS PÓS-TEXTUAIS
% -----------------------------------------------------------------------------
\postextual
\bibliography{referencias}

\begin{apendicesenv}
  \input{pos-textuais/apendice-a}
\end{apendicesenv}

\begin{anexosenv}
  \input{pos-textuais/anexo-a}
\end{anexosenv}

\end{document}
